% !PRINTED 418
\weberpage{418}
%<<CONT
and from this it follows, since $x$ and $x'$ are to be positive:
\[
x' = x, \quad y' = y, \quad (\eta = + 2).
\]
But if $-D = 4$, then (12) has in addition the solution
\[
\xi = \pm 1, \quad \eta = 0,
\]
and from (13) it follows:
\[
\begin{aligned}
y' &= \mp 2\,x, \\
x' &= \pm \frac{y}{2},
\end{aligned}
\]
and the sign of $\eta$ is to be determined so that $x'$ becomes positive. % [sic] The German print names \eta; since \eta=0 and x'=\pm y/2, it is actually the sign of \xi that determines the sign of x'.
If finally $-D = 3$, then besides (14) we have the solutions
\[
\xi = \pm 1, \quad \eta = \pm 1,
\]
in which at first both signs are arbitrary.
But then from (13) it follows:
\[
\begin{aligned}
\pm 2\,x' &= x \pm y, \\
\pm 2\,y' &= y \mp 3\,x,
\end{aligned}
\]
and one can choose the signs in two ways so that $x'$ becomes positive. We thus have the theorem:

I. \gesperrt{For a given equation (3), every proper solution of the equation (8) leads to a transformation $\left(\begin{smallmatrix}a, & b\\ c, & \partial\end{smallmatrix}\right)$ of the $n$th order. Different solutions of (8) lead in general to non-equivalent systems $\left(\begin{smallmatrix}a, & b\\ c, & \partial\end{smallmatrix}\right)$. In the two exceptional cases $D = -4$ and $D = -3$, in each case two or three different solutions of (8) lead to equivalent systems.}

If we denote by $k$ a number which is in general equal to the number of proper solutions of (8), but for $D = -4$ and $D = -3$ is half or a third of the number of these solutions, then we can state Theorem I. as follows:

II. \gesperrt{From an equation (3) we can derive $k$ non-equivalent systems $\left(\begin{smallmatrix}a, & b\\ c, & \partial\end{smallmatrix}\right)$ which satisfy the equation (1).}

\begin{center}
\textbf{\S\,115. The singular values of the invariant $j(\omega)$.}
\end{center}

The question that first claims our interest is that of the values of the modular functions of $\omega$ for the
%CONT>>
